\documentclass{article}
\usepackage[utf8]{inputenc}
\usepackage{graphicx}
\usepackage{hyperref}

\title{Documentazione del Progetto di Ingegneria del Software}
\author{Barichello Rebecca (VR471455) \\ Kaur Taniya (VR481429) \\ Magrograssi Federico (VR472578)}
\date{}

\begin{document}

\maketitle

\section*{Capitolo 1: Note introduttive}

Il sistema software in questione è stato creato per la gestione dei servizi erogati da una Casa della Salute, come visite mediche, prelievi e medicazioni. Si suppone che il software venga distribuito ai pazienti e al personale della Casa interessati, e che sia usufruibile tramite appositi terminali presenti nella struttura.

Gli utenti, per poter usufruire dei servizi, dovranno essere in possesso delle loro credenziali. Queste verranno generate dal personale di segreteria e comunicate agli utenti stessi.

\subsection*{Obiettivi del sistema}
Il sistema è progettato per gestire le attività di una Casa della Salute, dove operano più medici di base, supportati da infermieri per medicazioni e prelievi, e un ufficio di segreteria comune. La struttura ospita:
\begin{itemize}
    \item 3 ambulatori per visite ad adulti,
    \item 1 ambulatorio pediatrico,
\end{itemize}

\subsection*{Funzionalità principali}
Il sistema permette di:
\begin{itemize}
    \item Registrare e gestire i dati di medici, infermieri e pazienti.
    \item Gestire le prenotazioni per visite.
    \item Mantenere una storia delle visite e delle prestazioni effettuate.
    \item Fornire resoconti annuali delle prestazioni erogate ai pazienti.
\end{itemize}

\subsection*{Autenticazione e accesso}
Il sistema permette l'autenticazione ai pazienti tramite le credenziali fornite dalla segreteria. Il personale medico e infermieristico, invece, può accedere tramite le proprie credenziali personali, sempre nella sezione login.

\section*{Capitolo 2: Scenari d'uso}

In questo capitolo vengono descritti gli scenari d'uso del sistema software per ciascun attore coinvolto: infermiere, medico, paziente e segreteria. Si assume che ogni utente si sia autenticato correttamente prima di eseguire qualsiasi operazione.

\subsection*{2.1 Scenari d'uso da parte del personale infermieristico}

Il personale infermieristico può eseguire le seguenti operazioni:
\begin{itemize}
    \item \textbf{Autenticazione}: L'infermiere accede al sistema utilizzando le proprie credenziali.
    \item \textbf{Visualizzazione delle prestazioni da eseguire}: L'infermiere può visualizzare l'elenco delle prestazioni (prelievi e medicazioni) che deve eseguire.
    \item \textbf{Inserimento esito delle prestazioni}: Al termine di una prestazione, l'infermiere inserisce l'esito e le eventuali note nel sistema.
\end{itemize}

\subsection*{2.2 Scenari d'uso da parte del personale medico}

Il personale medico può eseguire le seguenti operazioni:
\begin{itemize}
    \item \textbf{Autenticazione}: Il medico accede al sistema utilizzando le proprie credenziali.
    \item \textbf{Inserimento esito della visita}: Dopo aver effettuato una visita, il medico inserisce l'esito e i dettagli della visita nel sistema.
    \item \textbf{Notifica esito visita}: Il sistema invia una notifica via e-mail al paziente per informarlo che l'esito della visita è disponibile.
    \item \textbf{Visualizzazione visite in programma}: Il medico può visualizzare l'elenco delle visite programmate per i propri pazienti.
    \item \textbf{Visualizzazione dati pazienti}: Il medico può accedere ai dati dei pazienti che ha in carico.
\end{itemize}

\subsection*{2.3 Scenari d'uso da parte del paziente}

Il paziente (o il suo referente adulto) può eseguire le seguenti operazioni:
\begin{itemize}
    \item \textbf{Autenticazione}: Il paziente accede al sistema utilizzando le proprie credenziali.
    \item \textbf{Prenotazione visita}: Il paziente può prenotare una visita con il proprio medico curante o con un sostituto, se il medico curante non è disponibile.
    \item \textbf{Prenotazione prelievi/medicazioni}: Il paziente può prenotare prelievi o medicazioni presso la struttura.
    \item \textbf{Visualizzazione esito visita}: Il paziente può visualizzare l'esito di una visita effettuata.
    \item \textbf{Storico visite}: Il paziente può accedere alla storia delle visite e delle prestazioni ricevute.
\end{itemize}

\subsection*{2.4 Scenari d'uso da parte del personale di segreteria.}

Il personale di segreteria può eseguire le seguenti operazioni:
\begin{itemize}
    \item \textbf{Autenticazione}: La segreteria accede al sistema utilizzando le proprie credenziali.
    \item \textbf{Gestione dati pazienti}: La segreteria può inserire, modificare o eliminare i dati dei pazienti.
    \item \textbf{Gestione prenotazioni visite}: La segreteria può gestire le prenotazioni delle visite, inclusa la modifica o la cancellazione.
    \item \textbf{Gestione dati medici e infermieri}: La segreteria può inserire, modificare o eliminare i dati del personale medico e infermieristico.
    \item \textbf{Visualizzazione visite programmate}: La segreteria può visualizzare l'elenco delle visite programmate nella struttura.
\end{itemize}

\end{document}